% ==================================================================
% id:       der-005
% topic:    נגזרות
% subtopic: הגדרת הנגזרת וגזירות
% solution: full
% ==================================================================
\begin{enumerate}
\item
האם לפונקציה $\ds{f(x) =\case{cc}{ x^{\frac{5}{3}} \sin\lr{\frac{1}{x}} & x \neq 0 \\ 0 & x = 0 }}$ יש נגזרת רציפה בנקודה \(x = 0\)?
    
\hint{חשבו את $f'(0)$ לפי ההגדרה ואת $f'(x)$ ל-$x\neq0$ לפי כללי הגזירה, ובדקו אם $\lim_{x\to0}f'(x)=f'(0)$; בדקו כל מחובר של $f'(x)$ בנפרד.}
\item מהם הערכים של פרמטר
    $\alpha\in\mathbb{N}$
    כך שעבורם נגזרתה של
    $$\ds{f(x) =\case{cc}{ x^{\alpha} \sin\lr{\frac{1}{x}} & x \neq 0 \\ 0 & x = 0 }}$$
    רציפה בנקודה
    $x=0$?
\hint{בדקו קודם לאילו $\alpha$ קיימת $f'(0)$ לפי ההגדרה, ואז לאילו מהם $f'(x)\to f'(0)$. השוו את החזקות של $x$ שמופיעות.}
\end{enumerate}

\ifsol
\begin{solbox}
\begin{enumerate}
\item
ראשית, נמצא מהי
$f'(x)$.
ל-
$x\neq0$
מתקיים
    $$.f'(x)= \frac{5}{3}x^{\frac{2}{3}}\sin(\tfrac{1}{x})-\frac{1}{x^{\frac{1}{3}}}\cos(\tfrac{1}{x})$$
נחשב את הנגזרת ב-
$x=0$
לפי ההגדרה:
\begin{align*}
    f'(0)&=\ \lim_{x\to 0}\frac{f(x)-f(0)}{x}\\
    &=\ \lim_{x\to0}\frac{f(x)}{x}\\
    &=\ \lim_{x\to0}x^{\frac{2}{3}} \sin\lr{\frac{1}{x}}\\
    &=0,
\end{align*}
כמכפלה של פונקציה אפסה בפונקציה חסומה.
אז קיבלנו כי
$$.\ds{f'(x) =\case{cc}{ \frac{5}{3}x^{\frac{2}{3}}\sin(\tfrac{1}{x})-\frac{1}{x^{\frac{1}{3}}}\cos(\tfrac{1}{x}) & x \neq 0 \\ 0 & x = 0 }}$$
עתה מטרתינו לברר האם
$f'$
רציפה ב-
$x=0$,
כלומר,
האם
$$.\lim_{x\to 0}f'(x)=f'(0)=0$$
הפונקציה
$f'$
מיוצגת כסכום של שתי פונקציות. לגבי הראשונה,
$$,h(x)=\frac{5}{3}x^{\frac{2}{3}}\sin\left(\frac{1}{x}\right)$$
די ברור כי גבולה ב-
$x=0$
הוא
$0$
(מדוע?),
לגבי השניה,
נסמנה
$g(x)=\frac{1}{x^{\frac{1}{3}}}\cos(\tfrac{1}{x})$,
יש חששות. נראה כי גבולה ב-
$x=0$
הוא
לא
$0$
(ביתר הדיוק, הגבול ב-
$x=0$
לא קיים כלל, אך כדי להפריך רציפות ב-
$x=0$
די להראות כי גבולה ב-
$x=0$
שונה מערכה ב-
$x=0$
שהוא, כפי שחישבנו,
$0$).
תהא
$a_n=\frac{1}{2\pi n}$,
עבורה
$\displaystyle\lim_{n\to\infty}a_n=0$
מתקיים
$$,\lim_{n\to\infty}g(a_n)=\lim_{n\to\infty}2\pi n\cos(2\pi n)=\lim_{n\to\infty}2\pi n=\infty$$
ולכן, ע"פ היינה,
$\displaystyle\lim_{x\to 0}g(x)$
לא קיים.
מכך נובע כי
$$\lim_{x\to 0}f'(a_n)=\lim_{x\to0}h(x)+\lim_{x\to0}g(x)$$
אינו קיים
(כי
$\displaystyle\lim_{x\to0}h(x)=0$).
מכאן נסיק כי
$f'$
לא רציפה ב-
$x=0$.

    \item
   ל-
   $x\neq 0$
   מתקיים
   $$.f'(x)= \alpha x^{\alpha-1}\sin(\tfrac{1}{x})-x^{\alpha-2}\cos(\tfrac{1}{x})$$
נחשב את הנגזרת ב-
$x=0$:
   \begin{align*}
    f'(0)&=\ \lim_{x\to 0}\frac{f(x)-f(0)}{x}\\
    &=\ \lim_{x\to0}\frac{f(x)}{x}\\
    &=\ \lim_{x\to0}x^{\alpha-1} \sin\lr{\frac{1}{x}}.
\end{align*}
נשים לב כי הגבול האחרון קיים
(ושווה ל-
$0$)
אמ"מ
$\alpha>1$,
ואמנם, אם
$\alpha>1$
הפונקציה היא מכפלה של פונקציה אפסה
$x^{\alpha-1}$
בפונקציה חסומה
$\sin(\tfrac{1}{x})$,
ולכן גבולה הוא
$0$,
ואם
$\alpha\leq1$
אין ל-
$x^{\alpha-1} \sin\lr{\frac{1}{x}}$
גבול ב-
$x=0$
(מדוע?).
אז קיבלנו כי
$$.\ds{f'(x) =\case{cc}{ \alpha x^{\alpha-1}\sin(\tfrac{1}{x})-x^{\alpha-2}\cos(\tfrac{1}{x}) & x \neq 0 \\ 0 & x = 0 }}$$
אנו כבר יודעים כי לגזירותה של
$f$
ב-
$x=0$
נדרש
$\alpha>1$.
עבור ערכים אלה של הפרמטר
$$.\lim_{x\to0}\alpha x^{\alpha-1}\sin(\tfrac{1}{x})=0$$
נשים לב כי
$$\lim_{x\to0}x^{\alpha-2}\cos(\tfrac{1}{x})$$
קיים ושווה ל-
$0$
(שהוא הערך של הנגזרת ב-
$x=0$)
אמ"מ
$\alpha>2$
(מדוע?).
נסיק כי
$f'$
רציפה ב-
$x=0$
עבור כל
$\alpha>2$.
\end{enumerate}
\end{solbox}
\fi
